\section{The Tower} \label{app:tower}

This appendix records the Cayley-Dickson construction used throughout, the single move that builds the whole algebra
tower and that the pinch-point argument rests on.

\subsection*{The doubling}

Given an algebra $A$ with a conjugation $x \mapsto \bar{x}$, the Cayley-Dickson double is the set of pairs
$(a, b)$ with $a, b \in A$, with multiplication
\begin{equation}
  (a, b)(c, d) = (ac - \bar{d}\,b, \ d a + b \bar{c}),
\end{equation}
and conjugation $\overline{(a,b)} = (\bar{a}, -b)$. The double has twice the dimension of $A$. Applied repeatedly from
the one-dimensional start it gives the tower
\begin{equation}
  \mathbb{R} \to \mathbb{C} \to \mathbb{H} \to \mathbb{O} \to \mathbb{S}, \qquad 1, 2, 4, 8, 16,
\end{equation}
the reals, complexes, quaternions, octonions, and sedenions. At the base we read $\mathbb{R}$ discretely as the
integers $\mathbb{Z}$, since the construction is discrete throughout and the continuum is emergent.

\subsection*{What each step loses}

Each doubling trades away one property, in a fixed order.

\begin{center}
\begin{tabular}{lll}
  \hline
  Step & Dimension & Property lost \\
  \hline
  $\mathbb{Z} \to \mathbb{C}$ & $1 \to 2$ & order (gains $i^2 = -1$) \\
  $\mathbb{C} \to \mathbb{H}$ & $2 \to 4$ & commutativity \\
  $\mathbb{H} \to \mathbb{O}$ & $4 \to 8$ & associativity \\
  $\mathbb{O} \to \mathbb{S}$ & $8 \to 16$ & the norm (composition) \\
  \hline
\end{tabular}
\end{center}

The order is forced by one fact. The double of $A$ keeps the norm-multiplicative property if and only if $A$ is
associative (the Hurwitz condition). The reals, complexes, and quaternions are associative, so the first three
doublings keep the norm, producing $\mathbb{C}$, $\mathbb{H}$, $\mathbb{O}$. The octonions are not associative, so the
fourth doubling loses the norm, and the sedenions have zero divisors. The composition algebras are therefore exactly
$\mathbb{R}, \mathbb{C}, \mathbb{H}, \mathbb{O}$, of dimensions $1, 2, 4, 8$. This is the ceiling at eight.

\subsection*{The octonion multiplication table}

For reference, the multiplication of the seven imaginary units $e_1, \dots, e_7$ used throughout the fermion
construction is recorded below. Each unit squares to $-1$, and the off-diagonal products follow the Fano-plane
orientation. The entry in row $e_i$, column $e_j$ is the product $e_i e_j$.

\begin{center}
\renewcommand{\arraystretch}{1.2}
\begin{tabular}{c|ccccccc}
        & $e_1$  & $e_2$  & $e_3$  & $e_4$  & $e_5$  & $e_6$  & $e_7$  \\
  \hline
  $e_1$ & $-1$   & $e_3$  & $-e_2$ & $e_5$  & $-e_4$ & $-e_7$ & $e_6$  \\
  $e_2$ & $-e_3$ & $-1$   & $e_1$  & $e_6$  & $e_7$  & $-e_4$ & $-e_5$ \\
  $e_3$ & $e_2$  & $-e_1$ & $-1$   & $e_7$  & $-e_6$ & $e_5$  & $-e_4$ \\
  $e_4$ & $-e_5$ & $-e_6$ & $-e_7$ & $-1$   & $e_1$  & $e_2$  & $e_3$  \\
  $e_5$ & $e_4$  & $-e_7$ & $e_6$  & $-e_1$ & $-1$   & $-e_3$ & $e_2$  \\
  $e_6$ & $e_7$  & $e_4$  & $-e_5$ & $-e_2$ & $e_3$  & $-1$   & $-e_1$ \\
  $e_7$ & $-e_6$ & $e_5$  & $e_4$  & $-e_3$ & $-e_2$ & $e_1$  & $-1$   \\
\end{tabular}
\end{center}

The table is antisymmetric off the diagonal, $e_i e_j = - e_j e_i$ for $i \neq j$, recording non-commutativity, and a
direct check of any associator such as $(e_1 e_2) e_4$ against $e_1 (e_2 e_4)$ shows non-associativity. These two
facts, the anticommutation and the non-associativity, are the inputs to the Clifford and triality arguments
respectively. The left-multiplication matrices $L_a$ of Section~\ref{sec:fermions} are read column by column from this
table.
